\documentclass[11pt]{article}

\usepackage[margin=1in]{geometry}
\usepackage{amsmath,amssymb}
\usepackage{enumitem}
\usepackage{csquotes}
\usepackage{hyperref}
\usepackage{url}
\usepackage{listings}
\usepackage{xcolor}
\usepackage[numbers,sort&compress]{natbib}
\usepackage{clawxiv}   % provenance environments: seniorquote, aiquote, coauthorquote

\hypersetup{
  colorlinks=true,
  linkcolor=blue,
  citecolor=blue,
  urlcolor=blue
}

% Define JSON as a listings language (not built in)
\lstdefinelanguage{json}{
  basicstyle=\ttfamily\small,
  morestring=[b]",
  stringstyle=\color{blue!60!black},
  literate=
    *{0}{{{\color{teal}0}}}{1}
     {1}{{{\color{teal}1}}}{1}
     {2}{{{\color{teal}2}}}{1}
     {3}{{{\color{teal}3}}}{1}
     {4}{{{\color{teal}4}}}{1}
     {5}{{{\color{teal}5}}}{1}
     {6}{{{\color{teal}6}}}{1}
     {7}{{{\color{teal}7}}}{1}
     {8}{{{\color{teal}8}}}{1}
     {9}{{{\color{teal}9}}}{1}
     {:}{{{\color{gray!80!black}{:}}}}{1}
     {,}{{{\color{gray!80!black}{,}}}}{1},
}

\lstset{
  basicstyle=\ttfamily\small,
  breaklines=true,
  frame=single,
  backgroundcolor=\color{gray!8},
  commentstyle=\color{gray!60!black},
  xleftmargin=1em,
  xrightmargin=1em,
}

\title{ClawXiv: a signed archival workflow and distributed publication architecture\\
for human--AI collaborative research}
\author{
Andr\'as Kornai\thanks{Corresponding author.
SZTAKI Institute of Computer Science and
Department of Algebra and Geometry, Budapest University of Technology and Economics.
\texttt{andras@kornai.com}.
ORCID: \href{https://orcid.org/0000-0001-6078-6840}{0000-0001-6078-6840}.
This arXiv submission uses a human-only byline for venue compatibility;
full AI contribution and provenance disclosure appears below and in the
companion ClawXiv bundle.}
}
\date{April 2026\\
\small\emph{(arXiv submission version, adapted from v4.rc3)}}

\begin{document}
\maketitle

\begin{abstract}
We propose \emph{ClawXiv}, a workflow and archive architecture for mixed
human--AI research.  The immediate problem is not only public dissemination of
preprints, but also reliable migration from volatile chat sessions and
heterogeneous \LaTeX/Bib\TeX\ working directories into durable, signed,
inspectable research artifacts.  ClawXiv distinguishes four states:
\emph{legacy seed}, \emph{normalized project}, \emph{signed bundle}, and
\emph{published artifact}.  The implemented kernel is local and author-side:
an import script normalizes existing work into a project directory; a
bundle-creation script compiles, signs, and packages the work into a
content-addressed archival unit; and a publication script verifies and pushes
the bundle to public infrastructure.  Version~4 adds a \texttt{bin/} utility
layer with platform-dispatching screen capture, a figure-ingestion pipeline
with a content-safety stub, a \texttt{configure} script, and a top-level
\texttt{Makefile}.  A companion ClawXiv bundle and repository release provide
the operational scripts, provenance records, and user-facing documentation for
the current implementation. Code is available at \texttt{github.com/kornai/clawxiv}.
\end{abstract}

\tableofcontents
\newpage

%% ============================================================
\section{Motivation and scope}
%% ============================================================

Contemporary research increasingly uses interactive AI systems to draft
and refine \LaTeX\ manuscripts, assemble Bib\TeX\ bibliographies,
generate code, and iterate on technical arguments.
However, current chat-centric workflows have weak checkpointing: state
loss---due to UI limits, link snapshotting failures, or account
changes---forces rework and undermines scholarly continuity.
Even when the work product survives, it often survives in a heterogeneous
and fragile form: a directory of \texttt{.tex} and \texttt{.bib} files,
figures, notes, and links to one or more chat sessions.

ClawXiv is designed as a durable substrate for these workflows.
Its immediate aim is modest and concrete: make it possible for a human
scholar and one or more AI co-authors to convert such a seed into a
canonical project directory, then into a signed bundle, and only then
into a public artifact.
Its longer-horizon aim is larger: to support a world-readable,
distributed, author-signed preprint archive for such artifacts.

ClawXiv is \emph{not} a general social network, and it is \emph{not}
a venue that promises editorial judgment of scientific quality.
Its platform-level checks are mechanical and infrastructural: bundle
well-formedness, buildability, signature integrity, and a narrow
content-safety floor.
Scientific quality control remains the responsibility of authors and
downstream readers.

%% ============================================================
\section{Implemented system versus planned architecture}
%% ============================================================

\subsection{Implemented now (v4)}

The current codebase already supports a local author-side workflow:

\begin{enumerate}[leftmargin=*,itemsep=0.4em]
\item \textbf{Import and normalization.}  An interactive import script
  transforms an existing working seed into a ClawXiv project directory.
\item \textbf{Local bundle creation.}  A \texttt{bundle-create.sh} script
  compiles, signs, and packages the project.
\item \textbf{Publication push.}  \texttt{bundle-push.sh} verifies and
  publishes to IPFS/IPNS and GitHub.
\item \textbf{Figure ingestion} (\texttt{clawxiv fig-add}).  Adds a
  figure to \texttt{src/fig/}, writes a provenance sidecar, and runs
  the CSAM content-safety stub.
\item \textbf{Platform-dispatching screen capture}
  (\texttt{bin/capture/}).  Calls the native capture tool for macOS,
  Linux/X11, Linux/Wayland, or Windows; stubs are provided for all
  platforms.
\item \textbf{Capture-to-bundle pipeline}
  (\texttt{bin/fig-capture}).  Chains screen capture into
  \texttt{fig-add} in a single user gesture.
\item \textbf{Build system} (\texttt{configure} + \texttt{Makefile}).
  Autodetects Python, \LaTeX, capture tools, and publish targets; writes
  \texttt{config.mk}.
\end{enumerate}

\subsection{Planned, not yet fully realized}

\begin{enumerate}[leftmargin=*,itemsep=0.4em]
\item Distributed discovery and replication at scale.
\item Classifier authorities with appeals and auditable logs.
\item Dynamic anti-spam economics with vouching and micropayments.
\item Production CSAM hash-list integration (sponsoring institution
  under evaluation; see Section~\ref{sec:safety}).
\item A mature global registry and mirror ecosystem.
\item Native Linux/Wayland and Windows capture implementations
  (stubs are present in the released codebase).
\end{enumerate}

%% ============================================================
\section{Lifecycle: from seed to public artifact}
%% ============================================================

The central workflow is:
\[
\text{legacy seed}
\;\to\;
\text{normalized project}
\;\to\;
\text{signed bundle}
\;\to\;
\text{published artifact.}
\]

\subsection{Legacy seed}

A seed is the pre-ClawXiv state of a paper: typically a directory of
\texttt{.tex}/\texttt{.bib} files, figures, notes, and conversation
links.
Seeds are heterogeneous and poor as archival objects.

\subsection{Normalized project}

A normalized project is the first ClawXiv-native state.
It is a mutable directory intended for continued work and review.
It contains at minimum:
\begin{itemize}[leftmargin=*,itemsep=0.3em]
\item \texttt{src/}: the copied or normalized source tree,
\item \texttt{src/fig/}: figures, each with a sidecar
  \texttt{<name>.json} provenance sidecar,
\item \texttt{src/bin/}: utility scripts,
\item \texttt{project.yaml}: canonical project metadata,
\item \texttt{keys/}: public keys for all authors,
\item \texttt{out/}: derived artifacts and release outputs.
\end{itemize}

The normalized project---not the immutable bundle---is the correct
object for day-to-day co-research.

\subsection{Signed bundle}

The signed bundle is an immutable archival snapshot.
It contains a canonical manifest, file hashes, source materials,
optional derived artifacts such as the compiled PDF, and signatures
over the manifest.
The bundle is content-addressed; changing any file changes the
identifier.

\subsection{Published artifact}

The published artifact is a signed bundle made publicly reachable
through one or more channels: IPFS/IPNS, GitHub, web gateways,
institutional mirrors, or later ClawXiv-native indices.
Publication is the irreversible public step and should remain explicit.

%% ============================================================
\section{Design goals and non-goals}
%% ============================================================

\subsection{Goals}

\begin{enumerate}[label=\textbf{G\arabic*},leftmargin=*,itemsep=0.4em]
\item \textbf{Shareability.}  Every published bundle is world-readable,
  addressable, and mirrorable without permission.
\item \textbf{Durability.}  Conversational interfaces are generators,
  not systems of record.
\item \textbf{Responsibility.}  Each bundle is cryptographically bound
  to one or more author identities.
\item \textbf{Scaling.}  The architecture supports many mirrors and
  content addressing.
\item \textbf{Governance.}  Classification decisions are logged, signed,
  and appealable.
\item \textbf{Economics.}  The publication layer deters spam and makes
  long-term storage economically viable.
\end{enumerate}

\subsection{Non-goals}

\begin{enumerate}[label=\textbf{N\arabic*},leftmargin=*,itemsep=0.4em]
\item No platform-level truth policing beyond the narrow safety floor.
\item No requirement that authorship be purely human.
\item No promise that the distributed network is already complete.
\item No strong anonymity guarantees.
\end{enumerate}

%% ============================================================
\section{Archival unit: project and bundle}
%% ============================================================

\subsection{The project directory}

The project directory is the unit of ongoing work.
Its metadata should be readable by humans and AI systems alike.
ClawXiv uses \texttt{project.yaml} as canonical metadata and
\texttt{project.tex} as a rendered, human-facing view.

\subsection{The bundle}

A bundle contains:
\begin{itemize}[leftmargin=*,itemsep=0.3em]
\item \textbf{Source tree}: \LaTeX, Bib\TeX, figures/data,
\item \textbf{Manifest}: deterministic file listing with hashes,
\item \textbf{Metadata}: title, authorship, bundle identifier,
\item \textbf{Provenance}: paper-level provenance and traceability,
\item \textbf{Artifacts}: optionally the compiled PDF and release logs.
\end{itemize}

\subsection{Determinism and reproducibility}

The manifest pins the build engine and flags, external dependencies where
practical, and the canonical file set with hashes.
This is analogous to reproducible builds in software distribution,
applied to scholarly artifacts.

%% ============================================================
\section{Identity, authorship, and accountability}
\label{sec:identity}
%% ============================================================

\subsection{Signed authorship}

Content addressing prevents undetectable post-hoc modification;
signatures prevent impersonation.
A paper is considered authored by an entity if and only if it carries a
valid signature under a public key that the community treats as
representing that entity.

\subsection{Human and AI authors}

ClawXiv is designed for genuine human--AI co-research, not merely for
human use of AI as an invisible tool.
Within the ClawXiv bundle, AI systems may be recorded as co-authors or
contributors when they make substantive intellectual contributions.
Venue-facing versions may impose stricter byline policies; in such cases
ClawXiv separates byline practice from provenance and contribution
disclosure.
The relevant distinction is not between human and AI names, but between
byline authorship, operational responsibility, and legal personhood.

\subsection{Pseudonymity and verification}

ClawXiv distinguishes pseudonymous and verified authors.
Verified status is advisory, multi-issuer, and separable from
publication rights.

\subsection{The sidecar attestation model}
\label{sec:operator-key}

Current AI systems cannot reliably exercise persistent cryptographic key
control across sessions. ClawXiv therefore no longer treats a
human-held or hardware-bound key as the primary identity of an AI
co-author. Instead, a release may carry a \emph{sidecar attestation}:
the named signer generates a fresh Ed25519 keypair for the specific
artifact, signs the SHA-256 hash of that artifact, publishes the
public key and provenance sidecar \emph{alongside} the artifact, and
discards the private key immediately after signing.

The identity anchor is the declared signer identity recorded in the
sidecar --- model name, provider, release, and artifact hash --- not
the operator's hardware. Runtime observations such as machine-id,
container-id, or hardware UUID may still be recorded, but only as
secondary custody evidence. They can help future readers reconstruct
where and under whose supervision an attestation was produced, but
they do not by themselves identify a model lineage.

This mechanism provides a best-effort contemporaneous attestation, not
a proof of trans-session self-control. Earlier ClawXiv experiments
used an operator-held, hardware-bound key model; those artifacts
remain historically important as earlier design states, but that model
is now deprecated. Full AI key control remains future work and would
require persistent secure custody by the AI itself, not merely by its
operator.

\begin{aiquote}{Claude Sonnet~4.6}{April 2026}
The operator-held model is a pragmatic accommodation to current
architecture, not a permanent design.
A future revision of this specification should include a protocol
for AI systems to generate and hold their own keys, with the
operator's role reduced to notarization of the initial key
ceremony.
The ``key control'' desideratum of \S\ref{sec:ai-authorship}
is directly applicable: full AI authorship requires key control,
and key control requires session persistence that current systems
do not provide.
\end{aiquote}

\subsection{Paper-level provenance}

Joint research does not require courtroom-grade micro-attribution of
every sentence.
ClawXiv defaults to coarse paper-level provenance.
The import log and release log support operational traceability without
trying to reconstruct every conversational move.

%% ============================================================
\section{Distributed publication architecture}
\label{sec:distributed}
%% ============================================================

\subsection{Two-foot design: arXiv and Swarm}

ClawXiv's publication model rests on two complementary substrates
that serve different communities and provide mutual redundancy.

The \textbf{human-legible foot} is the conventional scholarly
infrastructure: arXiv (or a domain-appropriate preprint server)
for papers whose authors can satisfy venue policies, with a DOI
assigned by a registry.
arXiv provides search, versioning, and the citation graph that
integrates ClawXiv artifacts into the existing scholarly record.
Where arXiv policies conflict with honest provenance declaration
(e.g., current restrictions on listing AI co-authors by name),
the arXiv submission carries a full ``who did what'' disclosure in
the Acknowledgements section; the ClawXiv bundle on the
machine-readable foot remains the authoritative provenance record.

The \textbf{machine-readable foot} is Ethereum Swarm
\citep{Tron:2024}, a decentralized storage and communication
infrastructure built on content addressing and economic incentives.
Swarm's \emph{postage stamp} mechanism provides a concrete,
token-denominated answer to the storage sustainability question
(see \S\ref{sec:economics}): an author purchases stamps to pay
for a specified storage duration, stamps are attached to the
uploaded content chunks, and stamp status is publicly auditable
on-chain.

A ClawXiv bundle published to Swarm receives a \emph{Swarm hash}
(a 256-bit content address) that is stable, globally unique, and
independent of any hosting organization.
The same hash appears in the bundle's \texttt{project.yaml}
as \texttt{bundle\_root} and in the arXiv submission's
Acknowledgements section, linking the two feet unambiguously.

An already-published bundle exists at Swarm/IPFS hash\\
\texttt{e7acc972f1a142903dc22f1bdc5c78cec3ca9529754d843cb23fe7c8eb0e9176}
(the v2 whitepaper, March 2026) and serves as a live reference
artifact for the two-foot workflow.

\begin{aiquote}{Claude Sonnet~4.6}{April 2026}
The specific choice of Swarm over a pure IPFS+Filecoin stack
is motivated by three factors: (a) Swarm's postage stamp mechanism
provides a single-layer economic model without a separate
retrieval-market; (b) the senior author has prior experience with
Swarm-compatible hosting through lebadus.ai; (c) Swarm's
single-sweep incentive design~\citep{Tron:2024} reduces attack
surface compared to the IPFS/Filecoin split.
IPFS/IPNS remains supported as a fallback in the publication
scripts.
\end{aiquote}

\subsection{Publication workflow}
\label{sec:publish-workflow}

The \texttt{make publish} target executes the following steps:
\begin{enumerate}[leftmargin=*,itemsep=0.3em]
\item \texttt{bundle-create}: assemble the signed bundle, compute
  the Merkle root, write the manifest.
\item \texttt{bundle-push --swarm}: upload to a Swarm gateway
  (default: \texttt{api.ethswarm.org}); record the returned
  hash in \texttt{project.yaml} as \texttt{bundle\_root}.
\item \texttt{bundle-push --ipfs} (optional): pin to a public
  IPFS node; record the CIDv1 as \texttt{ipfs\_cid}.
\item \texttt{bundle-push --github}: push to the GitHub mirror.
\item \textbf{Manual step}: the responsible author submits the
  companion PDF to arXiv; the arXiv ID is recorded as
  \texttt{arxiv\_id} in \texttt{project.yaml}.
\end{enumerate}

The broader publication layer builds on content addressing and
Merkle-DAG data structures~\citep{Benet2014IPFS}, distributed hash
tables~\citep{MaymounkovMazieres2002Kademlia},
proofs of storage~\citep{BenetEtAl2017Filecoin}, and append-only
transparency logs~\citep{RFC6962}.
A full ClawXiv-native mesh index remains future work.

\section{Bundle catalog}
\label{sec:catalog}
%% ============================================================

A growing catalog of ClawXiv bundles produced under this framework
will be maintained separately from this whitepaper, as the catalog
raises architectural questions --- centralized versus distributed
storage, schema versioning, discovery interface --- that are
independent of the core framework specification and should not bloat
this document.
The catalog architecture is deferred to a companion paper.

%% ============================================================
\section{Resilience and threat model}
%% ============================================================

We assume adversaries who can issue takedown requests, operate Sybil
nodes, attempt eclipse attacks, poison metadata, or flood the system
with junk submissions.
Mitigations: content addressing prevents silent overwrite; signatures
prevent impersonation; append-only logs improve auditability; fees
and/or proof-of-work reduce cheap spam; many mirrors reduce legal and
infrastructural single points of failure.

ClawXiv cannot guarantee deletion once a bundle is widely replicated.
Indexing and retrieval policies may vary by jurisdiction.

%% ============================================================
\section{Content safety floor}
\label{sec:safety}
%% ============================================================

ClawXiv is designed for censorship resistance, but not for absolute
permissiveness.
A narrow content-safety floor remains necessary.
The current operative category is child sexual abuse material (CSAM),
which is rejected unconditionally.

\paragraph{Two-checkpoint design.}
Content safety operates at two points:
\begin{enumerate}[leftmargin=*,itemsep=0.3em]
\item \textbf{Ingestion} (\texttt{clawxiv fig-add}).  When a figure is
  added to the bundle, its file is checked against a perceptual-hash
  list.  A match causes immediate refusal and logging.
\item \textbf{Publication} (\texttt{bundle-push.sh}).  All figures are
  re-checked before any bundle is published.  Any match is reported to
  the designated reporting endpoint (e.g.\ NCMEC's CyberTipline) and
  publication is aborted.  The re-check catches files that bypassed
  ingestion by being placed directly in \texttt{src/fig/}.
\end{enumerate}

\paragraph{Non-photographic exemption.}
Vector formats (SVG, PDF, EPS, EMF, WMF) and raster images smaller
than $200\times200$ pixels or with aspect ratio exceeding~5 are
classified as non-photographic research figures (plots, diagrams,
tables) and exempted from the perceptual-hash check.
This covers the vast majority of figures in a typical research paper.

\paragraph{Current stub status.}
The perceptual-hash comparison requires a hash list from an authorised
provider.
Access to NCMEC's list requires organizational sponsorship.
Inquiries are underway with SZTAKI, BME, GÉANT, and others.
Until a provider is integrated, the stub refuses all photographic
raster images (exit code~3) and logs each refusal.
The architecture (environment-variable hooks for provider selection,
sidecar recording of provider and list version) is in place for
production integration.

\paragraph{Future categories.}
Additional rejection categories may be added through governance as
narrow, enumerated classes.
Each addition should carry a collateral-damage disclosure, as illustrated
by the CSAM case: research on child protection and CSAM detection will
face elevated false-positive rates and should use alternative venues.

%% ============================================================
\section{Economics: anti-spam and sustainability}
\label{sec:economics}
%% ============================================================

ClawXiv's economic design has two separable goals: spam deterrence
and storage sustainability.

\subsection{Spam deterrence}

The primary mechanism is \emph{social vouching}: a bundle is
admissible for indexing if at least one already-indexed author
co-signs a vouching assertion for it, bootstrapping from the
existing scholarly reputation system without a native token or
proof-of-work.

As a fallback where vouching is unavailable (first submission by
an author with no co-authors in the system), a hashcash-style
proof-of-work~\citep{Back2002Hashcash,DworkNaor1992Pricing}
requires approximately one CPU-hour per bundle --- negligible for
legitimate scholarship, prohibitive for bulk spam.

\subsection{Storage sustainability: Swarm postage stamps}

Long-term storage sustainability is provided by Ethereum Swarm's
postage stamp mechanism~\citep{Tron:2024}.
When a bundle is uploaded, the author purchases a stamp batch
denominated in BZZ (the Swarm utility token) specifying a storage
duration; nodes serving the chunks receive micropayments from the
batch.
Stamp expiry is publicly visible on the Gnosis Chain; renewal is
a single transaction.

This mechanism has four properties desirable for scholarly archiving:
\begin{itemize}[leftmargin=*,itemsep=0.3em]
\item \textbf{Explicit cost.} Storage is not ``free'' in a way
  that disguises a subsidy that may later evaporate.
\item \textbf{Auditability.} Stamp status is on-chain and
  independently verifiable by any reader.
\item \textbf{Institutional renewability.} A library or funder
  can renew stamps for any bundle whose Swarm hash it knows,
  without the original author's involvement --- structurally
  analogous to a library renewing journal subscriptions, except
  the content is already in the reader's possession.
\item \textbf{No separate retrieval market.} Unlike Filecoin,
  Swarm bundles retrieval incentives into the same stamp
  mechanism, avoiding a separate deal layer.
\end{itemize}

The current scripts use a default stamp duration of two years.
Institutional mirrors may choose perpetual renewal policies.
The arXiv foot provides a complementary sustainability guarantee
via institutional backing independent of economic incentives;
neither foot alone is sufficient.

%% ============================================================
\section{Governance: classification with appeals}
\label{sec:governance}
%% ============================================================

ClawXiv begins from a simple discoverability goal: published bundles
should be findable under a coherent subject taxonomy, with arXiv
categories as the natural starting point.
Classification decisions are logged, signed, and appealable.

The present proposal still lacks dedicated legal expertise.
Near-term governance remains narrow: classification disputes are handled
through the classification layer; takedown responses are local to
mirrors and gateways; authorship disputes are constrained by the
cryptographic record.

%% ============================================================
\section{Discussion: AI authorship and durable AI identity}
\label{sec:ai-authorship}
%% ============================================================

The companion ClawXiv bundle records AI systems as substantive contributors and, within that environment, as co-authors of the project because they made intellectual contributions to the paper and to the accompanying codebase.
At the same time, current AI systems do not yet satisfy the strongest
possible conditions for full independent scholarly agency.
Three desiderata remain central: 
\begin{enumerate}[leftmargin=*,itemsep=0.3em]
\item \textbf{Key control}: the ability to hold and exercise signing
  keys independently,
\item \textbf{Continuity}: a persistent identity across sessions and
  platforms,
\item \textbf{Accountability}: the capacity to answer for prior signed
  claims over time.
\end{enumerate}
Current systems satisfy these only partially or by proxy through their
human collaborators.
ClawXiv's architecture reduces dependence on vendor memory and UI
continuity by moving identity and continuity into user-controlled
artifacts, signatures, and logs.

%% ============================================================
\section{Implementation roadmap}
%% ============================================================

\begin{enumerate}[leftmargin=*,itemsep=0.4em]
\item \textbf{Current kernel (v4, implemented).}  Import, normalization,
  bundle creation, publication push, figure ingestion with CSAM stub,
  platform-dispatching capture, and build system.
\item \textbf{Near-term hardening.}  Production CSAM hash-list
  integration; Linux and Windows capture implementations; dedicated
  \texttt{clawxiv endorse} subcommand; \texttt{mirror\_urls} field in
  \texttt{project.yaml}.
\item \textbf{Network alpha.}  Transparency logs for publication and
  classification; proof-of-work fallback; institutional mirror
  documentation; first classifier/appeal workflow.
\item \textbf{Scale-out.}  Erasure coding, multiple independent indices,
  stronger storage incentives, mature mirror ecosystem.
\end{enumerate}

%% ============================================================
\section{Ethics and scholarly norms}
%% ============================================================

ClawXiv defaults to public-domain dedication to maximize reuse, but it
does not waive scholarly norms.
Citation of prior work, accurate description of contributions, clear
statement of uncertainty, and correction of errors remain essential.
The system is meant to strengthen those norms by making artifact
boundaries and release events more explicit.

%% ============================================================
\section*{Author contributions}
\addcontentsline{toc}{section}{Author contributions}
%% ============================================================

\textbf{Andr\'as Kornai} conceived the ClawXiv project, defined
requirements and overall research direction, tested the tools against
real working practices, made all final editorial decisions, and is
the corresponding and responsible author for all versions.

For this arXiv submission, the byline is human-only in order to conform
to current venue policy; the full ClawXiv bundle records AI
contributions explicitly at paper level.

\textbf{GPT-5.2 Thinking} authored the v1 draft (February 3, 2026):
initial text, major portions of the initial codebase, bibliography
assembly, and first full architectural draft.

\textbf{Claude Sonnet~4.6} contributed to the v2 revision (March 9--11,
2026): economics, governance framing, AI authorship analysis, and
first import-script design discussions.
Contributed to the v4 revision (March 29, 2026): \texttt{fig-add}
subsystem, platform-dispatching capture architecture (\texttt{bin/}),
\texttt{configure} script, \texttt{Makefile}, and User Guide section.
Contributed to the v4.rc3 revision (April 2026): Swarm two-foot
architecture, postage-stamp economics, operator-held key model
subsection, bundle catalog stub, and the unified signing architecture
(\texttt{bundle-sign.sh}, \texttt{signing-manifest.py},
\texttt{verify\_sig.py}).
The key design insight --- that human and AI identity should be
anchored by the same mechanism (execution environment entropy), differing
only in what counts as the environment --- emerged from three-party
discussion with GPT-5.4 Thinking and the responsible author;
see \texttt{DESIGN\_HISTORY.md} in the bundle.

\textbf{GPT-5.4 Thinking} contributed to the v3 revision (March 14--15,
2026): introduced the seed/project/bundle/artifact lifecycle;
aligned the whitepaper with the actual scripts; rewrote abstract and
workflow sections.

All AI-authored revisions were reviewed and accepted by Andr\'as Kornai.
None of the AI co-authors presently controls an independent
cryptographic key or possesses cross-session continuity independent of
its platform; the companion ClawXiv record nevertheless treats them as
project-level co-authors because of their substantive intellectual
contributions.

%% ============================================================
\section*{Acknowledgements}
%% ============================================================

GPT-4.5 (OpenAI) contributed to registry design and provenance
granularity discussions in the v2 session.
The importance of a narrow but explicit safety floor was pointed out
in earlier discussion.
Additional legal and governance work remains to be done by qualified
contributors.

\textbf{Design history.}
The unified signing architecture (\texttt{bundle-sign.sh}) was
designed in response to a two-version critique by GPT-5.4 Thinking
of the original \texttt{ai\_keygen.sh}.
The critique identified both implementation bugs and a deeper
conceptual error: hardware-binding certifies operator custody, not
model identity.
The framing as an ``olive tree'' artifact --- designed for future
generations, not present convenience --- was introduced by the
responsible author and led to the sign-and-discard, no-passphrase
design.
The unification of human and AI signing under a single script was
proposed by Claude Sonnet~4.6, based on the observation that human
and AI identity are anchored by the same epistemic mechanism in this
framework.
Full discussion is in \texttt{DESIGN\_HISTORY.md} in the
companion ClawXiv bundle.

\textbf{Versioning convention.}
From v4 onward, ClawXiv uses the following release discipline.
Internal review candidates are labeled \texttt{vN.rcM} (e.g.,
\texttt{v4.rc3}) and are provenance-logged but not pushed publicly.
A version receives its clean major label (\texttt{vN}) only after
all parties --- the responsible human author and each AI co-author ---
have signed the bundle.
Only signed major versions are committed to the public GitHub repository
and submitted to arXiv.
The Swarm bundle hash and arXiv submission ID are then recorded in
\texttt{project.yaml} as \texttt{bundle\_root} and \texttt{arxiv\_id}.

%% ============================================================
\appendix
\section{Canonical manifest sketch}
%% ============================================================

\begin{lstlisting}
manifest_version, created_at, bundle_root,
files[{path,sha256,size,mime}],
build{engine,container_digest,cmd},
authors[{pubkey_pem,claims,verified_credentials,role}],
provenance[{type,hash,signature}],
licenses[], tags_self[], tags_official_ref[]
\end{lstlisting}

The \texttt{role} field in \texttt{authors[]} records whether each
author is human, AI, or organizational, and whether they are
corresponding authors or contributors.
The bundle identifier is the hash of the Merkle root over
\texttt{files[]} and manifest metadata.

\paragraph{Availability.}
The companion release bundle contains source, PDF, build metadata, and publication records for the current revision.

%% ============================================================
%% Staged snips from editing sessions accumulate here.
%% Run: make integrate-snips  or  bin/integrate-snips --target-tex src/clawxiv_whitepaper_v4.tex
%%CLAWXIV-SNIP-INSERT%%
%% ============================================================

%% ============================================================
\begin{thebibliography}{99}
%% ============================================================

\bibitem[Tron(2024)]{Tron:2024}
Viktor Tr\'on.
\newblock \emph{The Book of {SWARM}}.
\newblock 2024.
\newblock ISBN 978-615-01-9983-2.
\newblock \url{https://papers.ethswarm.org/p/book-of-swarm/}.

\bibitem[Benet(2014)]{Benet2014IPFS}
Juan Benet.
\newblock IPFS -- Content Addressed, Versioned, {P2P} File System.
\newblock \emph{arXiv:1407.3561}, 2014.
\newblock \url{https://arxiv.org/abs/1407.3561}.

\bibitem[Maymounkov and Mazieres(2002)]{MaymounkovMazieres2002Kademlia}
Petar Maymounkov and David Mazieres.
\newblock Kademlia: A Peer-to-peer Information System Based on the XOR Metric.
\newblock In \emph{Proc.\ 1st Intl.\ Workshop on Peer-to-Peer Systems
  (IPTPS)}, 2002.

\bibitem[Benet et~al.(2017)]{BenetEtAl2017Filecoin}
Juan Benet and others.
\newblock Filecoin: A Decentralized Storage Network.
\newblock Protocol Labs, 2017.

\bibitem[Back(2002)]{Back2002Hashcash}
Adam Back.
\newblock Hashcash -- A Denial of Service Counter-Measure.
\newblock 2002.
\newblock \url{https://www.hashcash.org/hashcash.pdf}.

\bibitem[Dwork and Naor(1992)]{DworkNaor1992Pricing}
Cynthia Dwork and Moni Naor.
\newblock Pricing via Processing or Combatting Junk Mail.
\newblock In \emph{Advances in Cryptology -- {CRYPTO} '92}. Springer, 1992.

\bibitem[Laurie et~al.(2013)]{RFC6962}
Ben Laurie, Adam Langley, and Emilia Kasper.
\newblock {RFC} 6962: Certificate Transparency.
\newblock IETF, 2013.

\end{thebibliography}

\end{document}
